\documentclass[12pt]{article}
\usepackage[T1]{fontenc}
\usepackage[utf8]{inputenc}
\usepackage{lmodern}
\usepackage{textcomp}
\usepackage{enumitem}
\usepackage{geometry}
\geometry{margin=1in}
\begin{document}
\begin{center}
Answer Key - Philosophy - Undergraduate Level Assessment 8 \
\small (Answers are ordered exactly as in the question paper.)
\end{center}

\section*{Multiple Choice}
\begin{enumerate}[label=\arabic*.]
\item \textbf{Answer:} D
\\ \textit{Options:} A) Fasting on Yom Kippur; B) Lighting Shabbat candles; C) The rainbow; D) Circumcision; E) The Torah
\\ \textit{Note:} Circumcision is the sign of the covenant for Jewish males as it symbolizes their unique relationship with God, following the biblical command given to Abraham as a physical manifestation of their faith and identity.
\item \textbf{Answer:} C
\\ \textit{Options:} A) useless.; B) valid.; C) ad hominem.; D) sound.
\\ \textit{Note:} Hare's assertion that all moral arguments are ad hominem highlights his belief that moral judgments are inherently tied to the character and motives of the individuals making those judgments, rather than being based solely on objective criteria.
\item \textbf{Answer:} A
\\ \textit{Options:} A) all actions are morally equivalent.; B) it is permissible to harm as a foreseen but unintended consequence of action.; C) the ends justify the means.; D) individuals should always act in their own best interest.; E) the moral worth of an action is determined by its outcome.
\\ \textit{Note:} Nagel argues that the absolutist position holds that all actions, regardless of context or consequences, are judged by the same moral standards, leading to the conclusion that they are morally equivalent.
\item \textbf{Answer:} A
\\ \textit{Options:} A) conditional and complex.; B) absolute and fundamental.; C) direct and secondary.; D) negotiable and fluctuating.; E) indirect but fundamental.
\\ \textit{Note:} Carruthers argues that our moral obligations to animals are not absolute but rather depend on various factors, including the context of human-animal relationships and the capacities of the animals involved, making them conditional and complex.
\item \textbf{Answer:} E
\\ \textit{Options:} A) our actions are completely predetermined; B) we have no control over our actions; C) the thesis of determinism is false; D) the thesis of free will is false; E) the thesis of determinism is true
\\ \textit{Note:} Soft determinism posits that while determinism is true and all events, including human actions, are determined by prior states of the world, individuals can still possess free will in a meaningful sense, thus affirming the truth of determinism.
\end{enumerate}

\section*{True / False}
\begin{enumerate}[label=\arabic*.]
\setcounter{enumi}{5}
\item \textbf{Answer:} False
\\ \textit{Options:} A) True; B) False
\\ \textit{Note:} The statement is false because Mani, the founder of Manichaeism, emphasized a dualistic cosmology rather than identifying God specifically as the Grand Architect focused solely on creation and design.
\item \textbf{Answer:} True
\\ \textit{Options:} A) True; B) False
\\ \textit{Note:} The statement is true because in predicate logic, 'Tdc' is a valid representation where 'T' denotes the teaching relation, 'd' represents David, and 'c' represents Chris, accurately capturing the relationship described in the original sentence.
\item \textbf{Answer:} True
\\ \textit{Options:} A) True; B) False
\\ \textit{Note:} Guru Nanak is recognized as the founder of Sikhism and the first of the ten Sikh Gurus because he established the core beliefs and practices of the faith in the 15 th century, laying the foundation for Sikh theology and community.
\item \textbf{Answer:} False
\\ \textit{Options:} A) True; B) False
\\ \textit{Note:} Hume argues that reason plays a role in moral decision-making, but he emphasizes that it is our emotions and sentiments that ultimately guide our moral choices, suggesting that reason and passion are interconnected rather than one being irrelevant.
\item \textbf{Answer:} False
\\ \textit{Options:} A) True; B) False
\\ \textit{Note:} The imposition of the death penalty does not guarantee the prevention of future murders because empirical evidence shows that it does not deter crime more effectively than other forms of punishment, and factors contributing to criminal behavior are complex and multifaceted.
\end{enumerate}

\section*{Long Form Response}
\begin{enumerate}[label=\arabic*.]
\setcounter{enumi}{10}
\item \textbf{Answer:} Hill's focus on the benefits of environmental preservation shifts the moral discourse from destruction to proactive stewardship by emphasizing the intrinsic and extrinsic advantages of maintaining ecological systems. By framing preservation as a beneficial choice, Hill encourages individuals and societies to recognize their responsibility in fostering a sustainable environment rather than viewing it merely as a resource to exploit. This perspective fosters a sense of ethical obligation that prioritizes long-term environmental health, which in turn supports biodiversity, enhances quality of life, and promotes intergenerational equity. Ultimately, this proactive stewardship cultivates a moral landscape where the health of the planet is seen as integral to human well-being, thus motivating collective action toward sustainable practices.
\\ \textit{Note:} correct because it articulates how Hill reframes environmental preservation as a moral imperative by highlighting its benefits, thereby promoting a shift from exploitative attitudes to a proactive, responsible approach to stewardship.
\item \textbf{Answer:} To analyze the validity of the argument represented by H ≡ (I · \textasciitilde{}J), K ≡ \textasciitilde{}H, and \textasciitilde{}(K ⊃ J) using indirect truth tables, we begin by assuming that I ≡ J is false, meaning I and J differ in truth values. If I is true and J is false, H becomes true, K is false, and \textasciitilde{}(K ⊃ J) holds true. However, if I is false and J is true, H becomes false, K is true, leading to a true implication in K ⊃ J, contradicting our assumption. In both cases, we encounter inconsistencies that prevent the premises from being true while allowing I and J to differ. Hence, the argument is valid, as it consistently leads to a conclusion that I must indeed be equivalent to J, leaving no room for counterexamples.
\\ \textit{Note:} The answer correctly demonstrates the validity of the argument by using indirect truth tables to show that assuming I and J have different truth values leads to contradictions with the given premises, thereby confirming that I ≡ J must hold true.
\end{enumerate}

\vfill
\noindent\textit{End of Answer Key}
\end{document}